% SOURCE_AUTHORITY: sga5_fr_workpass.tex, lines 4092--4127 (LF-normalized unit).
% SOURCE_SHA256: 791F4EFFC5E02832D5D77ED03518C8156D6F07E4C8238B03545DB93D883FBB28
\subsection*{6.13}
Supongamos que $S$ es el espectro de una clausura algebraica del cuerpo finito $\mathbb F_q$ y que $X/S$ procede, por extensión de escalares, de un esquema propio y liso $X_0/\mathbb F_q$. Tomemos como $F$ el endomorfismo de Frobenius de $X/S$ ``definido por la elevación de las coordenadas a la potencia $q$-\'esima''. Se tiene $X^F=X_0(\mathbb F_q)$ y $X^F$ está formado por puntos transversales, puesto que $dF=0$. Si $L$ es un haz localmente libre de tipo finito sobre $X$ y $u\in\Hom(F^*L,L)$, la fórmula de Woods Hole se escribe, por tanto,
\begin{equation}
\tag{6.13.1}
\sum_i(-1)^i\operatorname{Tr}((F,u),H^i(X,L))=
\sum_{x\in X_0(\mathbb F_q)}\operatorname{Tr}(u_x,L_x).
\end{equation}
En particular, para $L=\OO_X$ y $u$ el isomorfismo canónico $F^*\OO_X\xrightarrow{\sim}\OO_X$, se obtiene:
\begin{equation}
\tag{6.13.2}
\sum_i(-1)^i\operatorname{Tr}(F,H^i(X,\OO_X))=\operatorname{Card}(X_0(\mathbb F_q))\quad(\operatorname{mod} p).
\end{equation}
La teoría de Artin--Schreier permite reescribir el primer miembro como $\sum_i(-1)^i\operatorname{Tr}(F,H^i(X,\mathbb F_p))$, donde $H^*(X,\mathbb F_p)$ designa la cohomología étale de $X$ con valores en el haz constante $\mathbb F_p$. Mediante un refinamiento de este argumento, Deligne, en (SGA $4^{1/2}$ Fonctions L), dedujo de 6.13.1 la siguiente generalización de 6.13.2:


\begin{statement}{Teorema 6.13.3 (Deligne)}
Sean $X_0$ un esquema separado y de tipo finito sobre $\mathbb F_q$, $X$ el esquema deducido de $X_0$ por extensión de escalares a una clausura algebraica $k$ de $\mathbb F_q$ y $F$ el endomorfismo de Frobenius de $X/k$. Entonces, para todo $L_0\in\ob D_{\ctf}(X_0,\mathbb F_p)$, denotando por $L$ la imagen inversa de $L_0$ sobre $X$, se tiene
\[
\sum_i(-1)^i\operatorname{Tr}(F,H_c^i(X,L))=
\sum_{x\in X^F}\operatorname{Tr}(F,L_x).
\]
\end{statement}

Señalemos que el caso particular de 6.13.3 correspondiente a $L_0=\mathbb F_p$ fue establecido por Katz en (SGA 7 XXII) mediante un método completamente distinto, basado en el cálculo, a la manera de Dwork, de la función zeta de una hipersuperficie. Por otra parte, 6.13.2, en el caso propio y liso, es también consecuencia inmediata de la fórmula de Lefschetz en cohomología cristalina ([5] VII 3.1.9). De hecho, la cohomología cristalina proporciona la información más precisa de que cada factor $\det(1-Ft,H^i(X_0/W))\in W[t]$ de la función zeta es congruente módulo $p$ con $\det(1-Ft,H^i(X_0,\OO_{X_0}))$ (al menos si $H^*(X_0/W)$ no tiene torsión). Katz demostró [11] que, mediante las estimaciones de B. Mazur [14], se podían obtener congruencias superiores (en las que intervienen los $H^i(X_0,\Omega^j_{X_0})$ y la operación de Cartier) para hipersuperficies (o intersecciones completas) lisas cuyos números de Hodge (en la dimensión de interés) se suponen nulos hasta cierto grado.

Por otra parte, si $X_0$ es un esquema propio sobre $\mathbb F_q$, geométricamente íntegro y tal que $H^i(X_0,\OO_{X_0})=0$ para $i>0$, la fórmula de Lefschetz--Verdier 6.7.2, aplicada a $X=X_0\otimes k$ y al endomorfismo de Frobenius de $X$, implica inmediatamente que $X_0$ posee al menos un punto racional sobre $\mathbb F_q$ (puesto que $\operatorname{Tr}(F,R\Gamma(X,\OO_X))=1$). Serre pregunta si esta conclusión sigue siendo cierta cuando se sustituye $\mathbb F_q$ por un cuerpo que satisface $(C_1)$ ([17] II 3), por ejemplo $\mathbb C(t)$ (Tsen) o $\mathbb C((t))$ (Lang) (véase loc. cit. para otros ejemplos); observa, en efecto, que, si $X_0$ es una intersección completa de multigrado $(m_1,\ldots,m_r)$ en $\mathbb P^n$, la condición de trivialidad de la cohomología equivale ([16] n\textsuperscript{o} 78) a la condición $m_1+\cdots+m_r<n+1$, que es también la que interviene cuando se escribe la condición $(C_1)$.

\subsection*{6.14}
La fórmula de Woods Hole tiene muchas otras aplicaciones, para las cuales remitimos a [2] y a la exposición de Beauville [4].


\subsection*{6.15}

El redactor no ha examinado el caso de las componentes de puntos fijos de dimensión $\geq1$, pero no se excluye que de 6.7.2 (combinada probablemente con Riemann--Roch) puedan deducirse fórmulas de Lefschetz--Riemann--Roch al estilo de las de Donovan [7], Iversen [10] y Nielsen [15] (al menos con valores en el cuerpo base). Por otra parte, no se excluye que la fórmula 6.7.2, aplicada sobre los números duales, proporcione fórmulas de residuos para los campos vectoriales análogas a las de Bott [6], [3].

